\begin{question}

  \begin{questionpart}
    Prove that $(1/\sqrt{2})(1+i\sigma_x)$, where the matrix $\sigma_x$
    is given in (3.50), acting on a two-component spinor can be
    regarded as the matrix representation of the rotation operator
    about the $x$-axis by angle $\pi/2$. (The minus sign signifies
    that the rotation is clockwise.)
  \end{questionpart}

  \begin{questionpart}
    Construct the matrix representation of $S_z$ when the eigenkets of
    $S_y$ are used as base vectors.
  \end{questionpart}

\end{question}

\begin{purpose}
  Spin operators have a rule of $720^{\circ}$ in which they have to
  go around twice to come back to where they started.  Yeah, I know,
  it's like when we say that there is less than $360^{\circ}$ around
  the earth because of spatial warping associated with general
  relativity.  Some things are best appreciated when you start viewing
  the math.  This and other quirks below.  The trouble with this
  problem is that it really belongs in chapter 3.  We get a glimpse of
  some of these things, but chapter 1 really doesn't give us the tools
  we need to properly analyze the systems under consideration.
\end{purpose}

\begin{answer}

  \begin{answerpart}{The World is Spinoring}
    And recalling equation 3.50 from the book:
    
    \begin{equation*}
      \tag{3.50}
      \sigma_1 =
      \begin{pmatrix}
        0 & 1 \\
        1 & 0 \\
      \end{pmatrix},
      \sigma_2 =
      \begin{pmatrix}
        0 & -i \\
        i &  0 \\
      \end{pmatrix},
      \sigma_3 =
      \begin{pmatrix}
        1 &  0 \\
        0 & -1 \\
      \end{pmatrix}
    \end{equation*}

    Recalling equation \ref{1.11:spinAngle} and looking at figure
    \ref{1.11:spinAxes} we can rotate around the $x$ axis by setting
    $\beta$ to $-\pi/2$ and setting $\alpha$ to $-\pi/2$ so that our
    rotation is about the $x$ axis and a negative value for $\beta$
    results in a clockwise rotation.

    \begin{equation}
      \begin{split}
        \ket{\vectorbold{S}\cdot\vectorbold{\hat{n}}; +}
        &= \cos{\frac{\beta}{2}} \ket{+} + e^{i\alpha}\sin{\frac{\beta}{2}} \ket{-}
        \\
        &= \cos{-\frac{\pi}{4}}\ket{+} - i\sin{-\frac{\pi}{4}} \ket{-}
        \\
        &= \frac{1}{\sqrt{2}}\ket{+} + i\frac{1}{\sqrt{2}}\ket{-}
        \\
        &\dot{=} \frac{1}{\sqrt{2}}
        \begin{bmatrix}
          1\\i\\
        \end{bmatrix}
        \\
      \end{split}
    \end{equation}

    Applying the same equation for a rotation of $\pi+\pi/2=3\pi/2$ to
    see what we get with $\ket{-}$

    \begin{equation}
      \begin{split}
        \ket{\vectorbold{S}\cdot\vectorbold{\hat{n}}; -}
        &= \cos{\frac{\beta}{2}} \ket{+} + e^{i\alpha}\sin{\frac{\beta}{2}} \ket{-}
        \\
        &= \cos{-3\frac{\pi}{4}} \ket{+} - i\sin{-3\frac{\pi}{4}} \ket{-}
        \\
        &= -\frac{1}{\sqrt{2}}\ket{+} + i\frac{1}{\sqrt{2}}\ket{-}
        \\
        &\dot{=} \frac{1}{\sqrt{2}}
        \begin{bmatrix}
          -1\\i\\
        \end{bmatrix}
        \\
      \end{split}
    \end{equation}

    So we're looking for our operator $O$ to behave as follows

    \begin{equation}
      \begin{alignedat}{2}
        &&\hat{O}
        \begin{bmatrix}
          1 & 0\\ 0 & 1\\
        \end{bmatrix}
        &=
        \frac{1}{\sqrt{2}}
        \begin{bmatrix}
          1 & -1\\ i & i\\
        \end{bmatrix}
        \\
        \implies&& \hat{O}
        &=
        \frac{1}{\sqrt{2}}
        \begin{bmatrix}
          1 & -1\\ i & i\\
        \end{bmatrix}
        \\
      \end{alignedat}
    \end{equation}

    Trying what the they gave us in the problem

    \begin{equation}
      \begin{split}
        \hat{O}
        &= 
        \frac{1}{\sqrt{2}}
        \left(
        1 + i
        \begin{bmatrix}
          0 & 1 \\
          1 & 0 \\
        \end{bmatrix}
        \right)
        \\
        &= 
        \frac{1}{\sqrt{2}}
        \left(
        \begin{bmatrix}
          1 & 0 \\
          0 & 1 \\
        \end{bmatrix}
        +
        \begin{bmatrix}
          0 & i \\
          i & 0 \\
        \end{bmatrix}
        \right)
        \\
        &= 
        \frac{1}{\sqrt{2}}
        \begin{bmatrix}
          1 & i \\
          i & 1 \\
        \end{bmatrix}
      \end{split}
    \end{equation}

    Clearly...something has gone horribly wrong.  Let's take a look
    back at our solution from problem \ref{1.11:theQuestion}.  We
    started out assuming that we could put in any angle we wanted for
    $\beta$ (and $\alpha$ for that matter).  But consider this, the
    point in that diagram described by $\beta=\pi/2$ and
    $\alpha=\pi/2$ (which is the $y$ axis) is equivalent to
    $\beta=-\pi/2$ and $\alpha=-\pi/2$.  Same point, but anywhere we
    have a $\sin$ in a formulation, it's going to end up with a
    negative value owing to the odd nature of the function.  So,
    instead of targeting $-3\pi/2$, we're instead going to target
    $\pi/2$ which is geometrically equivalent, but reamins within the
    $[0,\pi]$ window for which the formula was developed.
    Interestingly, I don't actually see any specific line in the
    derivation from problem \ref{1.11:theQuestion} that depends upon
    this assumption of range, but the sign issue is there nonetheless.
    Let's see what happens when we correct for that by going $\pi/2$
    in the other direction instead.

    \begin{equation}
      \begin{split}
        \ket{\vectorbold{S}\cdot\vectorbold{\hat{n}}; +}
        &= \cos{\frac{\beta}{2}} \ket{+} + e^{i\alpha}\sin{\frac{\beta}{2}} \ket{-}
        \\
        &= \cos{\frac{\pi}{4}}\ket{+} + (\cos(-\frac{\pi}{2}) + i\sin(-\frac{\pi}{2}))\sin{\frac{\pi}{4}} \ket{-}
        \\
        &= \frac{1}{\sqrt{2}}\ket{+} - i\frac{1}{\sqrt{2}}\ket{-}
        \\
        &\dot{=} \frac{1}{\sqrt{2}}
        \begin{bmatrix}
          1\\-i\\
        \end{bmatrix}
        \\
      \end{split}
    \end{equation}

    OK...that still doesn't look right.  The whole result is off by a
    factor of $i$.  Going back to our foundational assumption

    \begin{equation}
      \inlabel{theAssumption}
      \ket{\vectorbold{S}\cdot\vectorbold{\hat{n}}; +}
      =
      \cos(\frac \beta 2)\ket + + \sin(\frac \beta 2) e^{i\alpha}\ket -
    \end{equation}

    we see that it's nooot quite right...  Let's see what happens when
    we use this formula and the geometry provided by the problem to
    construct $\ket{-}$ from $\ket{+}$.  It doesn't have the ambiguity
    we discussed above.

    \begin{equation}
      \begin{split}
        \ket{-}
        &= \cos(\frac{\beta}{2})\ket{+}
        +  e^{i\alpha}\sin(\frac{\beta}{2})\ket{-}
        \\
        &= \cos(\frac{\pi}{2})\ket{+}
        +  e^{0i}\sin(\frac{\pi}{2})\ket{-}
        \\
        &= 0\ket{+}
        +  1\ket{-}
        \\
        &= \ket{-}
        \\
      \end{split}
    \end{equation}

    So far so good right?  Watch what happens when we try to do that
    exact same rotation, but first we rotate about the $z$ axis before
    doing so.
    
    \begin{equation}
      \begin{split}
        \ket{-}
        &= \cos(\frac{\beta}{2})\ket{+}
        +  e^{i\alpha}\sin(\frac{\beta}{2})\ket{-}
        \\
        &= \cos(\frac{\pi}{2})\ket{+}
        +  e^{i\pi/2}\sin(\frac{\pi}{2})\ket{-}
        \\
        &= (\cos(\pi/2) + i\sin(\pi/2))\ket{-}
        \\
        &= i\sin(\pi/2)\ket{-}
        \\
        &= i\ket{-}
        \\
      \end{split}
    \end{equation}

    Where did that factor of $i$ come from?  This turns out to be
    pretty closely linked with classical mechanics when we talk about
    the consequences of going from point $A$ to point $B$ by different
    rotations.  The meaning and cause of this factor is pretty subtle
    and really isn't explored in chapter 1.  To be honest, this
    problem \emph{does not} belong in chapter 1.  When we go back to
    equation 1.102 in \S 1.4.2 of the book

    \begin{equation}
      \tag{1.102} \ket{S_x;+} = \frac{1}{\sqrt{2}}\ket{+} +
      \frac{1}{\sqrt{2}}e^{i\delta_1}\ket{-}
    \end{equation}

    We find that random phasor in there of $e^{i\delta_1}$.  He also
    adds in one for $\delta_2$ later.  Right after that equation, he
    states that ``In writing (1.102) we have used the fact that the
    \emph{overall} phase (common to both $\ket{+}$ and $\ket{-}$) of a
    state ket is immaterial; the coefficient of $\ket{+}$ can be
    chosen to be real and positive by convention.''  This particular
    section always felt a bit hand-wavey and that's probably because
    it doesn't really get developed until chapter 3 (where this
    problem belongs).  So, we can easily draw two conclusions here: 1)
    there's something semi-arbitrary about the selection and 2) it's
    really only useful for \emph{relative} phase and that there is
    some kind of \emph{global} phase that is ignored in this
    formulation.

    So what can we do?  Well, there are two paths forward that I can
    see.  One is we can bring in chapter 3 and just abbreviate
    equation 3.62, plug in $\phi$ and call it a day

    \begin{equation}
      \tag{3.62.abbrev}
      \exp(\frac{-i\sigma\cdot\hat{\vb{n}}\phi}{2})
      = \vb{1}\cos(\frac{\phi}{2})
      - i\sigma\cdot\hat{\vb{n}}\sin(\frac{\phi}{2})
    \end{equation}

    ooooor we could do it the hard way...or the really hard way.
    Let's just do it the hard way.  When we applied our formula above,
    we actually got our operator just right with the exception that we
    have a factor of $-i$ for the $\ket{-}$ term.  That factor is a
    consequence of doing our rotation to $\ket{-}$ through the $yz$
    plane rather than the $xz$ plane.  So if we correct for that
    phase, we end up with exactly the right answer.  Let's first see
    if our phase factor from the extra $\pi$ radians of rotation
    really is the culprit.  We took $\alpha$ to be $-\pi/2$ so that
    any negative values of $\beta$ would result in clockwise rotations
    so let's see what we're introducing by going past $\ket{-}$

    \begin{equation}
      \begin{split}
        \ket{-}
        &= \cos(\frac{\beta}{2})\ket{+}
        +  e^{i\alpha}\sin(\frac{\beta}{2})\ket{-}
        \\
        &= \cos(\frac{\pi}{2})\ket{+}
        +  e^{-i\pi/2}\sin(\frac{\pi}{2})\ket{-}
        \\
        &= (\cos(-\pi/2) + i\sin(-\pi/2))\ket{-}
        \\
        &= -i\sin(\pi/2)\ket{-}
        \\
        &= -i\ket{-}
        \\
      \end{split}
    \end{equation}

    So, when we applied our formulation earlier, we inadvetently
    introduced an additional factor of $-i$.  So, let's try putting
    together our two adjustements: the $\beta=\pi/2$ instead of
    $-3\pi/2$ and adding in a factor $-1/i = i$ to offset what we
    introduced by going around $\ket{-}$

    \begin{equation}
      \begin{split}
        \hat{O}\ket{-}
        &= i
        \left(
        \cos{\frac{\beta}{2}} \ket{+} + e^{i\alpha}\sin{\frac{\beta}{2}} \ket{-}
        \right)
        \\
        &= i
        \left(\cos{\frac{\pi}{4}} \ket{+} + \left(\cos(-\frac{\pi}{2}) + i\sin(-\frac{\pi}{2})\right)\sin{\frac{\pi}{4}} \ket{-}
        \right)
        \\
        &= \frac{i}{\sqrt{2}}
        \left(\ket{+} - i\ket{-}\right)
        \\
        &= \frac{1}{\sqrt{2}}
        \left(i\ket{+} + \ket{-}\right)
        \\
        &\dot{=} \frac{1}{\sqrt{2}}
        \begin{bmatrix}
          i\\1\\
        \end{bmatrix}
        \\
      \end{split}
    \end{equation}

    Substituting that into our earlier expression, we find that

    \begin{equation}
      \begin{alignedat}{2}
        &&\hat{O}
        \begin{bmatrix}
          1 & 0\\ 0 & 1\\
        \end{bmatrix}
        &=
        \frac{1}{\sqrt{2}}
        \begin{bmatrix}
          1 & i\\ i & 1\\
        \end{bmatrix}
        \\
        \implies&& \hat{O}
        &=
        \frac{1}{\sqrt{2}}
        \begin{bmatrix}
          1 & i\\ i & 1\\
        \end{bmatrix}\qed
        \\
      \end{alignedat}
    \end{equation}

    Is this answer a bit hand-wavey... yeppers.  I just don't see a
    better way to do this problem using chapter 1's material.

  \end{answerpart}

  \begin{answerpart}{\texorpdfstring{Building $Z$'s from $Y$'s}{Building Z's from Y's}}
    So... yeah... this problem doesn't really seem to provide much
    value but let's do it anyway.

    \begin{equation}
      \begin{alignedat}{2}
        &&\ket{y\pm} &= \frac{1}{\sqrt{2}} [\ket + \pm i\ket -]\\
        &&\ket{y+} + \ket{y-} &= \frac{1}{\sqrt{2}} [2\ket{+}]\\
        \implies&&\ket{+} &= \frac{\sqrt{2}}{2}[\ket{y+} + \ket{y-}]\\
        \implies&&\bra{+} &= \frac{\sqrt{2}}{2}[\bra{y+} + \bra{y-}]\\
        &&\ket{y+} - \ket{y-} &= \frac{1}{\sqrt{2}} [2i\ket{-}]\\
        \implies&&\ket{-} &= -i\frac{\sqrt{2}}{2}[\ket{y+} - \ket{y-}]\\
        \implies&&\bra{-} &= i\frac{\sqrt{2}}{2}[\bra{y+} - \bra{y-}]\\
      \end{alignedat}
    \end{equation}

    \begin{equation}
      \begin{split}
        \frac{2S_z}{\hbar}
        &= \Lambda_+ - \Lambda_-\\
        &= \dyad{+}{+} - \dyad{-}{-}\\
        &= \left(\frac{\sqrt{2}}{2}\right)^2
        \left[
          (\ket{y+} + \ket{y-})(\bra{y+} + \bra{y-})
          -
          (-i)(i)(\ket{y+} - \ket{y-})(\bra{y+} - \bra{y-})
          \right]
        \\
        &= \frac{1}{2}
        \left[
          (\ket{y+} + \ket{y-})(\bra{y+} + \bra{y-})
          -
          (\ket{y+} - \ket{y-})(\bra{y+} - \bra{y-})
          \right]
        \\
        &= \frac{1}{2}
        \left[
          (\cancel{\Lambda_{y+}} + \dyad{y-}{y+} + \dyad{y+}{y-} + \cancel{\Lambda_{y-}})
          -
          (\cancel{\Lambda_{y+}} - \dyad{y+}{y-} - \dyad{y-}{y+} + \cancel{\Lambda_{y-}})
          \right]
        \\
        &= \cancel{\frac{1}{2}}
        \left[
          \cancel{2}\dyad{y-}{y+}
          + \cancel{2}\dyad{y+}{y-}
          \right]
        \\
        &= \dyad{y-}{y+} + \dyad{y+}{y-}
        \\
      \end{split}
    \end{equation}

    Referencing equation \ref{1.7:theMatrixHasYouNow} to build our
    matrix, we have:

    \begin{equation}
      \begin{split}
        S_z
        &\dot{=}
        \begin{bmatrix}
          \mel{y+}{S_z}{y+} & \mel{y+}{S_z}{y-}\\
          \mel{y-}{S_z}{y+} & \mel{y-}{S_z}{y-}\\
        \end{bmatrix}
        \\
        &=
        \frac{\hbar}{2}
        \begin{bmatrix}
          0 & 1\\
          1 & 0\\
        \end{bmatrix}
        \\
        &=
        \frac{\hbar}{2}\sigma_1\qed
      \end{split}
    \end{equation}

    So, what he's trying to show us here is the cyclic nature of the
    spin operators.  We've just recreated $\sigma_x$ or $\sigma_1$.  I
    believe this will all become clear when we dig into the Bloch
    sphere, but until then I think he just wants you to look at that
    and say ``oooh shiny''.  Pretty sure Pauli spin operators were in
    chapter 1 in a previous incarnation and he didn't have a chance to
    renormalize before he died.  Once again, the world is a lesser
    place for his premature passing.
  \end{answerpart}
  
\end{answer}
